\section{Conclusions}
\begin{frame}{Security must follow the entire context-to-edit workflow}
\begin{columns}[T]
\column{.48\textwidth}
\textbf{For IDE providers}
\begin{itemize}
\item Make context sources and exclusions visible.
\item Filter suspicious context and suggestions.
\item Flag security-sensitive changes before acceptance.
\item Evaluate multi-step edits and incomplete remediation.
\end{itemize}
\column{.48\textwidth}
\textbf{For developers and teams}
\begin{itemize}
\item Treat suggested edits as unverified code.
\item Review the full transaction and skipped sites.
\item Check secrets, authorization, and migration settings explicitly.
\item Combine review, security analysis, and secure-coding training.
\end{itemize}
\end{columns}
\source{Section 6; transaction-level review follows the reported failure mechanisms}
\note{These recommendations are implications and proposed defenses, not defenses experimentally validated by this paper.}
\end{frame}

\begin{frame}{Three messages to take away}
\begin{enumerate}
\item \textbf{NES broadens the security boundary.} Ambient interactions and project context shape proactive edits.
\item \textbf{Better base models do not resolve every workflow risk.} Context propagation and lost security invariants appear across the tested systems.
\item \textbf{Review must cover the sequence.} Fluent code, navigation, or silence does not establish secure completion.
\end{enumerate}
\bigskip
\centering\Large\color{Teal}A Tab press accepts an edit;\\it does not establish its security.
\source{Sections 4--8}
\end{frame}
